\chapter{推荐资料}
\label{ch:material}
下面列出一些关于光学、电子学和通信的实用且易于入门的资源。

\section{光学}

\subsection*{Elements of Photonics}
Keigo Iizuka 教授撰写的这本教材系统介绍了实用光学器件背后的理论，覆盖面很广。尤其推荐其中关于掺铒光纤放大器（EDFA）增益与噪声特性的内容，因为 EDFA 是现代光学与光通信中的核心器件。


\subsection*{Modern Optics lectures by Prof. Partha Roy Chaudhuri}
这套 \href{https://www.youtube.com/watch?v=2WiMeh1Dxl8&list=PLbRMhDVUMngePMuAGeAUeGVuZffTFY-5i}{免费的 60 讲视频课程} 覆盖了麦克斯韦方程、偏振、双折射等基础内容，也讲到波导和电光调制器等更高级的话题。其结构清晰、数学推导扎实，非常值得推荐。

\subsection*{Non-linear Optics lectures by Prof. Samudra Roy}
这套 \href{https://www.youtube.com/watch?v=EiIDScj124Q&list=PLbRMhDVUMngfBwyonVP8VIsabtnsV3GVv}{60 讲视频课程} 先介绍线性光学中的基本概念，再从麦克斯韦方程出发详细推导二阶与三阶非线性效应，内容涵盖二次谐波产生（SHG）、差频产生（DFG）、光学参量振荡器（OPO）等本教程未涉及的现象。尤其推荐其中关于 $\chi^{(2)}$ 和 $\chi^{(3)}$ 非线性对称性与张量性质的讲解。

\subsection*{Nonlinear Fiber Optics}
Govind P. Agrawal 的这本经典教材对式~\ref{eq:GNLSE} 背后的数学做了极其细致的阐述，并频繁联系实验结果。若想深入了解偏振相关效应、受激布里渊散射以及新型光纤等高级主题，这是非常值得精读的一本书。

\subsection*{MIT "Demonstrations in Lasers and Optics" lecture series}
由 Shaoul Ezekiel 教授主讲的 \href{https://www.youtube.com/watch?v=1cEXNLP5uE0&list=PL4E7FAAD67B171EBC}{49 讲系列视频} 通过大量实际演示介绍自由空间光学中的基本现象，例如偏振、干涉和衍射。其系统化的实验展示非常值得参考。

\subsection*{International School on Parametric Non Linear Optics lectures}
这套 \href{https://www.youtube.com/@ispnlo9041/videos}{44 讲视频课程} 由多位非线性光学专家主讲，涵盖了许多普通课程中未必会系统介绍的高级主题。

\subsection*{Les' Lab}
这个 \href{https://www.youtube.com/@LesLaboratory/videos}{YouTube 频道} 主要做激光技术的动手演示，包括染料激光器、光谱仪、电光调制器、二次谐波产生，甚至超连续谱产生。其对相关电子系统的细致讲解也很有参考价值。

\subsection*{YourFavouriteTA}
这是本教程作者运营的 \href{https://www.youtube.com/@yourfavouriteta/videos}{YouTube 频道}，内容结合了实验演示、理论推导与数值模拟，旨在用直观方式解释通常只能在复杂方程中看到的线性与非线性光学过程。


\section{电子学}
\subsection*{The Signal Path}
这个 \href{https://www.youtube.com/@Thesignalpath/videos}{YouTube 频道} 专注于电源、信号发生器、探测器、示波器、频谱分析仪等实验室常见设备的实测、拆解与维修，尤其推荐其中与射频器件相关的视频。

\subsection*{EMPossible}
这个 \href{https://www.youtube.com/@empossible1577/playlists}{YouTube 频道} 既讲电磁理论基础，也讲射频波导、半导体能带结构、有限元分析等高级内容。它覆盖面广，图示清晰，教程风格紧凑流畅。

\subsection*{Keysight Electronics Tutorials}
Keysight 官方的 \href{https://www.youtube.com/@KeysightLabs/videos}{YouTube 频道} 重点讲解电信号测量。尤其推荐其中有关示波器的视频，因为示波器是可视化光电探测器测得“功率随时间变化”的常用工具。


\subsection*{Rohde\&Schwarz Electronics Tutorials}
Rohde\&Schwarz 提供的这套
\href{https://www.youtube.com/watch?v=rUDMo7hwihs&list=PLKxVoO5jUTlvsVtDcqrVn0ybqBVlLj2z8}{YouTube 播放列表} 深入讲解了电子器件和信号性质的测量与表征，例如功率、相位噪声、放大器噪声系数等。许多概念都可以直接迁移到光学领域。



\section{通信}


\subsection*{Ian Explains}
这个 \href{https://www.youtube.com/@iain_explains/videos}{YouTube 频道} 讲解数字与模拟信号处理、通信和统计学中的基础方法，例子直观、推导清楚，非常适合入门。


\section{仿真工具}

\subsection*{Octave Photonics}
这个免费的、交互式的 \href{https://www.octavephotonics.com/nlse}{网页工具} 可以直接求解式~\ref{eq:GNLSE}，非常适合用来快速可视化色散、孤子形成、拉曼效应等核心现象。

\subsection*{ssfm\_functions.py}
这是作者维护的 \href{https://github.com/OleKrarup123/NLSE-vector-solver/tree/main}{开源 Python 库}，可以灵活求解式~\ref{eq:GNLSE}。它支持通过编程构造复杂的光纤链路，例如多个跨度、不同拉曼模型、输入输出增益、色散补偿等，并带有方便的可视化函数。对于需要针对特定研究系统做建模的非线性光学研究生来说，这种可修改性尤其有价值。

\subsection*{gnlse-python}
这是一个 \href{https://github.com/WUST-FOG/gnlse-python}{开源 Python 库}，同样可用于求解式~\ref{eq:GNLSE}，具有非常好的 \href{https://gnlse.readthedocs.io/en/latest/index.html}{文档} 和内置可视化功能。它还有一篇介绍实现细节的 arXiv 预印本~\cite{redman2021gnlsepythonopensourcesoftware}。

\subsection*{GMMNLSE-Solver}
这个 \href{https://github.com/WiseLabAEP/GMMNLSE-Solver-FINAL}{开源 MATLAB 库} 可以求解式~\ref{eq:GNLSE} 的多模版本，从而建模多模孤子等更复杂的现象。它属于较为高级的工具，更适合已经非常熟悉本教程所讨论单模情形的读者。如果配置得当，还可以把计算并行到 GPU 上，对这类复杂仿真来说速度相当可观。


\section{其他}

\subsection{Essential Graduate Physics}
这本 \href{https://commons.library.stonybrook.edu/egp/}{免费教材} 系统讲解经典力学、电动力学、量子力学和统计物理。其图示质量高，计算过程清晰，即使是量子计算这类高级主题也写得相当易懂。
